% =====================================================================
%  notacion.tex — Notación común de la asignatura (todos los temas)
%  Procesamiento de Señales · Grado en Ing. Electrónica de Comunicaciones
% =====================================================================
%
%  Convenios (ver también comun/README.md):
%   * Tiempo continuo entre paréntesis x(t); tiempo discreto entre
%     corchetes x[n]. Nunca x(n) para una secuencia.
%   * Frecuencia angular continua Omega (rad/s) y frecuencia de
%     oscilación F (Hz); frecuencia angular discreta omega (rad/muestra)
%     y frecuencia normalizada f (ciclos/muestra). No se usa «w».
%   * CTFT X(jΩ), DTFT X(e^{jω}), Laplace X(s), Z X(z), siempre
%     bilaterales y con su región de convergencia (ROC).
%   * DFT: X[k] = Σ_{n=0}^{N-1} x[n] e^{-j2πkn/N}; IDFT con 1/N.
%     N es siempre una longitud; el orden de un filtro es M.
%   * La unidad imaginaria es j; la constante de Euler, e (recta).
%   * Un sistema se denota con T{·}: y(t) = T{x(t)}.
% ---------------------------------------------------------------------

% Conjuntos numéricos
\newcommand{\R}{\mathbb{R}}
\newcommand{\Z}{\mathbb{Z}}
\newcommand{\N}{\mathbb{N}}
\newcommand{\C}{\mathbb{C}}

% Constantes y diferenciales en redonda
\newcommand{\e}{\mathrm{e}}
\newcommand{\dd}{\mathop{}\!\mathrm{d}}
\newcommand{\dv}[2]{\frac{\mathrm{d}#1}{\mathrm{d}#2}}
\newcommand{\dvn}[3]{\frac{\mathrm{d}^{#3}#1}{\mathrm{d}#2^{#3}}}

% Parte real/imaginaria, argumento y funciones especiales
% (unicode-math redefine \Re e \Im al empezar el documento: se fijan después)
\AtBeginDocument{\renewcommand{\Re}{\operatorname{Re}}\renewcommand{\Im}{\operatorname{Im}}}
\DeclareMathOperator{\sinc}{sinc}
\DeclareMathOperator{\Sa}{Sa}
\DeclareMathOperator{\sgn}{sgn}
\DeclareMathOperator{\atantwo}{atan2}

% Delimitadores
\DeclarePairedDelimiter{\abs}{\lvert}{\rvert}
\DeclarePairedDelimiter{\norm}{\lVert}{\rVert}
\DeclarePairedDelimiter{\set}{\{}{\}}

% Convolución (lineal y circular)
\newcommand{\conv}{\ast}
\newcommand{\cconv}[1]{\circledast_{#1}}

% Sistemas
\newcommand{\T}{\mathcal{T}}
\newcommand{\sis}[1]{\T\set*{#1}}

% Transformadas y pares transformados
\newcommand{\Four}{\mathcal{F}}
\newcommand{\Lap}{\mathcal{L}}
\newcommand{\Zt}{\mathcal{Z}}
\newcommand{\DFT}{\operatorname{DFT}}
\newcommand{\IDFT}{\operatorname{IDFT}}
\newcommand{\ROC}{\mathrm{ROC}}
\newcommand{\parF}{\mathrel{\overset{\mathcal{F}}{\longleftrightarrow}}}
\newcommand{\parL}{\mathrel{\overset{\mathcal{L}}{\longleftrightarrow}}}
\newcommand{\parZ}{\mathrel{\overset{\mathcal{Z}}{\longleftrightarrow}}}
\newcommand{\parSF}{\mathrel{\overset{\mathcal{SF}}{\longleftrightarrow}}}

% Transformadas de Fourier
\newcommand{\Xc}{X(j\Omega)}           % CTFT
\newcommand{\Xd}{X(\e^{j\omega})}      % DTFT
\newcommand{\ejw}{\e^{j\omega}}
\newcommand{\ejW}{\e^{j\Omega}}

% Señales básicas
\newcommand{\escalon}{u}
\newcommand{\rampa}{r}
\newcommand{\rect}{\Pi}
\newcommand{\tri}{\Lambda}

% Unidades frecuentes (texto)
\newcommand{\Hz}{\,\mathrm{Hz}}
\newcommand{\kHz}{\,\mathrm{kHz}}
\newcommand{\MHz}{\,\mathrm{MHz}}
\newcommand{\GHz}{\,\mathrm{GHz}}
\newcommand{\rads}{\,\mathrm{rad/s}}
